\chapter{Discussion}

\FloatBarrier

\section{Main findings and contributions}

\subsection{Graph-aware prompting and label recovery}

RQ1 receives a consistent answer: graph-aware prompting improves all five main metrics in every model-by-style comparison, including all three Llama seeds. Each effect uses matched completed responses, so the two arms of a contrast contain the same entities. This supports the versioned prompting and evaluation design in C1 and C3, while the empty-child control shows why the effect should be described as a combination of structural evidence and instructions.


The experiment supports the use of a taxonomy as an inference-time evidence source. Parents can supply a semantic boundary and useful broader vocabulary. Siblings and children can describe nearby distinctions, while combined context allows the model to relate those distinctions to the description. The source of the improvement may include label availability, reduced uncertainty about the intended branch, and better ordering. The protocol does not isolate these mechanisms, so it would be premature to describe the gain as proof of graph reasoning in an internal cognitive sense.


Oracle selection is central to this interpretation. Removing the centre label from the context prevents direct copying of that field, but leaves an informative neighbourhood chosen using gold. The resulting gains justify testing retrieval: there is useful evidence available if the right structure can be located. They do not yet show that a model can locate it reliably or remain accurate when the supplied centre is wrong.


\subsection{Prompt style and model differences}

For RQ4, Gemini has higher paired F1 than every Llama seed in each of the nine conditions. Gemini’s highest graph F1 occurs with an explanation-bearing template, while the standard graph template leads the three combined-graph styles for Llama at every seed. These patterns support the matched empirical comparisons in C2, not a general architectural explanation. Gemini’s standard outputs already involve internal thinking, so asking for a visible explanation is not equivalent to turning reasoning on.


\subsection{Relation evidence and depth}

For RQ2, parent-only F1 exceeds sibling-only F1 throughout, which does not support H1. Combined context is not uniformly better than parent-only across Llama seeds. The seed-0 child intervention has slightly negative F1 point estimates, but every interval includes zero. A precision–recall trade-off is observed for standard prompting only. For RQ3, gains generally shrink with depth, with the largest gain in the shallowest stratum in 10 of 12 comparisons; H2 is not supported.


\subsection{Interpreting benchmark agreement}

The 265 descriptions containing the normalized centre label and the 735 descriptions without it both receive explicit analysis. Sixteen listed gold sets are not ancestor-closed, and 89 centres differ from Gmax(e). These checks explain why context selection, flat recovery, and hierarchical evaluation must remain distinct. Deterministic scores make the comparison reproducible, while semantic review is needed before interpreting every unmatched label as false.

